The exponent \(\omega\) of matrix multiplication is the least real number
such that, for every \(\eps>0\), two \(n\times n\) matrices can be
multiplied with \(O(n^{\omega+\eps})\) arithmetic operations. A preprint
posted by OpenAI in October 2026 announced the bound \(\omega\le9/4\).
Its proof concludes through an existence argument for tensor characters
and does not specify an algorithm. Building on
the convolution relations and the shared-leg separator of that preprint,
we give a self-contained construction by explicit tensor maps: for every
rational \(\eps>0\), a terminating procedure outputs an exact bilinear
algorithm with rational coefficients whose exponent is below
\(9/4+\eps\). The construction uses two kinds of tensor relations among
products of polynomial multiplication tensors (splits and triplings).
How often each relation is applied is given by the expected visit counts
of a random walk whose transition probabilities come from the least
solution of the inequalities that these relations impose. These
frequencies are rounded to integers without breaking the balance of any
tensor factor. For every integer \(a\ge2^{64}\) we obtain a scheme that
computes several products of \(B\times B\) matrices from products of
smaller \(A\times A\) matrices; its recursive use has exponent below
\(9/4+900/\log a\) (natural logarithm), and
\(\log A\le181a^5\lceil\log_2a\rceil^2\). We give no bound on the size of
the final algorithm or on the time needed to construct it, and all
results concern exact arithmetic.
